% English translation of checked_pages/esq_119_029.tex; source: cote 119, site 29, PDF page 30, typescript page 27.
% AI-translated by GPT-6 Astra (OpenAI, ChatGPT Pro mode).
\scanpage{119/29}
\mspage{27}

(without boundary) Y, with, as additional structure linking them,
the so-called ``gluing'' map
\[ f : \partial U \longrightarrow Y \]
which is a proper smooth fibration. The initial situation
(X,Y) is reconstructed from $(U,Y, f : \partial U \longrightarrow Y)$ by the
formula
\[ X \simeq U \amalg_{\partial U} Y \]
(amalgamated sum under $\partial U$, mapping to U and Y through the inclusion,
resp. the gluing map).

This naive vision immediately encounters various
difficulties. The first is the somewhat vague nature of the
very notion of a tubular neighborhood, which acquires an at all
precise meaning only in the presence of much more rigid structures
than the topological structure alone, such as the ``piecewise
linear'' or riemannian structure (more generally, that of a space with a
distance function); the trouble here is that in none of the examples
that come spontaneously to mind does one naturally have
a structure of this kind - at best one has an equivalence class
of such structures, allowing one to rigidify the situation
just a little. If one moreover admits that an expedient
has been found for obtaining a tubular neighborhood with the desired
properties, which is also unique modulo an automorphism
(topological, say) of the situation, an automorphism that also
respects the fibered structure supplied by the gluing function,
there remains the difficulty that the choices made are not canonical,
the automorphism in question visibly failing to be unique, whatever
one does to ``normalize'' it. The idea here, to make canonical
what is not so, is to work systematically in
``\uline{isotopy categories}'' associated with the categories of a topological
kind that arise in these questions (such as the category of
admissible pairs (X,Y) and homeomorphisms of such pairs\frapparatus{V-FR-029-001},
etc.), keeping the same objects but taking as ``morphisms''

\margin{29}

